\documentclass{article}

\usepackage[T1]{fontenc}
\usepackage[french]{babel} 
\usepackage{mathptmx}
\usepackage{geometry}
\usepackage{graphicx}
\geometry{a4paper, margin=2cm}

\usepackage[
    colorlinks=true,        
    linkcolor=black,
    urlcolor=blue,          
    citecolor=green,        
    pdftitle={Rapport du Sprint de Décembre 2026}, 
]{hyperref}

\usepackage[acronym, toc]{glossaries}
\makeglossaries

\newacronym[
    description={Delay-Tolerant Networking. Réseau tolérant aux délais et aux ruptures de connectivité, où les nœuds ne sont pas forcément connectés en permanence. Les données sont transmises de manière opportuniste via des contacts planifiés.}
]{dtn}{DTN}{Delay-Tolerant Networking}

\newacronym[
    description={Contact Plan. Fichier de configuration décrivant les fenêtres de communication planifiées entre nœuds : qui communique avec qui, à quelle heure, pendant combien de temps, avec quel débit et quel délai.}
]{cp}{CP}{Contact Plan}

\newacronym[
    description={Schedule-Aware Bundle Routing. Algorithme de routage pour réseaux DTN qui utilise des métriques de distance (temps d'arrivée, nombre de hops) pour trouver les meilleures routes entre nœuds. \cite{sabr}}
]{sabr}{SABR}{Schedule-Aware Bundle Routing}

\newacronym[
    description={Structure de données centrale du projet représentant le réseau. Contient les nœuds (senders/receivers) et les contacts entre eux. Permet plusieurs contacts entre une même paire de nœuds à des moments différents, d'où le préfixe "multi".}
]{mg}{Multigraph}{Multigraph}

\newacronym[
    description={Mécanisme du Dynamic Loader des systèmes UNIX permettant de charger un module dont les fonctions seront utilisées en priorités avant celles de la glibC}
]{ldp}{PRELOAD}{LD\_PRELOAD}

\newacronym[
    description={Bundle Protocole Socket. Type spécifique de socket utilisée dans le cas des DTN}
]{bp}{BP Socket}{Bundle Protocole Socket}

\title{Contributions on Schedule-Aware Bundle Routing \\ \Large Projet de Programmation 2 }
\author{DAGUE Bastien, AL-FARAJ Samy, JOUINI Yassine, BROCHARD Vincent \\ \small L3 Informatique - Faculté des Sciences - Université de Montpellier}
\date{2025 - 2026}

\begin{document}

\maketitle
\tableofcontents

\newpage
\printglossary[type=\acronymtype, title={Glossaire et Abréviations}, style=altlist]

\newpage
\section{Introduction}
\subsection{A-SABR}
Notre projet de programmation porte sur la librairie ASABR, portée par Mr Olivier De Jonckère du LIRMM.

Le réseau de communication spacial repose sur le \gls{dtn}. Avec ce type de protocole, le routage s'effecture par des algorithmes de type \gls{sabr}.\\

La librairie A-SABR implémente donc ce type d'algorithme en proposant une grande modularité et une bonne performance.\cite{sabr}
C'est également pour suivre ce choix de performance que celle-ci est programmé en Rust, language connue pour son compilateur stricte et sa fine gestion de la mémoire.

\subsection{Sujet}

L'intérêt pour ce sujet a été immédiat au sein de notre groupe, motivé par la perspective de contribuer à un projet Open Source d'envergure réelle. 
Ce cadre de travail nous a offert l'opportunité de nous initier au langage Rust, 
dont la montée en puissance dans l'industrie logicielle souligne l'importance actuelle. 
Pour plusieurs membres de l'équipe aspirant à se spécialiser en informatique système (bas niveau) ou en 
réseaux, ce projet constituait une application concrète et exigeante de nos ambitions professionnelles.\\

Le sujet présentait plusieurs point à aborder : \\
\begin{itemize}
    \item Portage au no-std
    \item Divers tâche de DevOps (Tests, Pipeline)
    \item Création d'un nouvel onglet dans l'interface graphique DT-Chat
    \item Intégration dans le logiciel Hardy \\
\end{itemize}


Au final, tout à été un peu chamboulé, le no-std à été abbandonnée, l'intégration à Hardy jugée peut intéressante. \\

Nous avons surtout compris qu'ici on ne dépendait pas seulement de nous mais qu'il s'agissait d'un immense travail d'équipe.
On commence une tâche, on la trouve accompli, on lance une Pull Request, on discute, on modifie, tous ça prend un certain temps et nécessite beaucoup de vas et viens.\\

Arrivé à la fin du projet voici les réels objectifs accompli :  \\

\begin{description}
    \item[Interface graphique DT-Chat :] Développement d'un volet « Options » dans l'interface graphique DT-Chat : Ce module permet désormais de modifier dynamiquement, au cours de l'exécution, l'algorithme de routage utilisé ou le \gls{cp} chargé.
    
    \item[Industrialisation (CI/CD) :] M    Mise en place d'une pipeline CI/CD robuste pour assurer la pérennité de la bibliothèque. Ce travail a été complété par une campagne de tests unitaires couvrant les modules principaux du projet.\\

    \item[Expérimentation Système (BP Socket) :] 
    Pendant le dernier mois du projet, nous avons travaillé sur \gls{bp}.

    \gls{bp} telle qu'elle à été développée utilise un module kernel nécessitant 
    les droit root pour être utilisée.
    L'objectif est de permettre l'usage de \gls{bp} sur des systèmes embarqués contraints où les privilèges « root » ne sont pas disponibles. L'expérimentation du mécanisme \gls{ldp}
    vise à intercepter les appels système pour s'affranchir de la dépendance au module noyau actuel.

    Ce dernier objectif est une total découverte de ce mécanisme offert par le Dynamic Loader et 
    toute avancée même non complète intéresse Mr De Jonckère.
\end{description}


\newpage
\section{Développements Logiciel : Conception, Modélisation, Implémentation}
\subsection{DT-Chat}
\subsubsection[DT-Chat Backend]{Backend}
\subsubsection[DT-Chat Frontend]{Frontend}
\subsection{Pipeline}
Bastien : Ici detail point par point la pipeline l'évolution de ta PR, précise bien sur que ça à été PR etc etc

\subsection{Test A-SABR}
\subsubsection[Module : Contact Manager]{Contact Manager}
Ici ta partie Yassine, que tu partagera surement avec Bastien et/ou Samy mais c'est ton module

\subsubsection[Module : Pathfinding]{Pathfinding}
\paragraph{Try\_make\_hop :}~\\\\
Ici ta partie Samy

\paragraph{Contact,Hybrid et Node parenting :}~\\\\

\subsection{BP-Socket}

\begin{thebibliography}{9} 

\bibitem{sabr} 
J. A. Fraire, O. De Jonckère, and S. Burleigh \textit{Routing in the space internet: A
contact graph routing tutorial.}, Journal of Network and Computer Applications, 174, 102884.

\bibitem{sabr} 
O. De Jonckère, L. Ma, and J. A. Fraire. (1994) \textit{\LaTeX : -SABR: The Adaptive Library for Schedule-Aware Bundle
Routing.}, WiSEE 2025 - IEEE International Conference on Wireless for Space and Extreme Environments, Oct 2025,
Halifax, Canada.

\end{thebibliography}
\end{document}